\documentclass[12pt,a4paper]{article}
\usepackage{amsmath,amsfonts,amssymb}
\usepackage{graphicx}
\usepackage{hyperref}
\usepackage{enumitem}
\usepackage{geometry}
\geometry{margin=2.5cm}
\emergencystretch=2em

\begin{document}

\begin{center}
    {\Large \textbf{Master Thesis Project in Computational Physics}}\\[0.3cm]
    {\large \textbf{Physics-Informed Neural Networks for the Time Evolution of Interacting Electrons in Two-Dimensional Multi-Well Quantum Dots}}\\[0.2cm]
    \vspace{0.3cm}
    {Duration: One Year (60 ECTS)}
\end{center}

\vspace{1cm}

\section*{1. Introduction}

Electrons confined in electrostatically defined quantum dots, and in
particular electrons trapped above the surface of superfluid helium,
are candidates for quantum computing hardware~[1--3].  In the designs
studied in our group each electron sits in its own potential well, the
two lowest states of its in-plane motion define a qubit, and the
unscreened Coulomb interaction between electrons in neighboring wells
provides the coupling.  Gates are carried out by changing the shape of
the wells in time.  Predicting how such a device behaves requires the
solution of the time-dependent Schr\"odinger equation for several
interacting electrons in a potential that changes in time.

Two lines of work in the group address this problem from different
sides.  The first is the multi-configuration time-dependent Hartree
and Hartree--Fock method, MCTDH(F)~[4--6], developed for multi-well
devices in the PhD work of Oskar Leinonen~[3,7].  It treats a register
of four wells, with one or three electrons per well, including the
full time dependence of the gates.  The electrons move in one
dimension, with the transverse motion integrated out into an effective
interaction, and the cost grows quickly with the number of electrons
per well and with the dimension of the grid.  The second line is the
neural-network approach of Ref.~[8], which combines a
Slater--Jastrow--backflow wave function with a physics-informed
training stage and has produced accurate ground states for up to
twenty electrons in a single two-dimensional dot.  It works directly
in two dimensions with the bare Coulomb interaction and has no grid,
but it has so far been used only for time-independent problems in one
well.

This thesis project brings the two together.  The aim is to extend the
physics-informed neural-network (PINN) calculations of Ref.~[8] to
electrons in two dimensions with one electron in each well, as in
Ref.~[7], to use them for the time evolution of two and then more
electrons under time-independent and time-dependent external probes,
and to compare the results with those of the MCTDH(F) codes.

\section*{2. Physical System and Scientific Motivation}

With $\hbar=m_e=1$, and with electron $\sigma$ confined to well
$\sigma$ centered at $\boldsymbol r_c^{\sigma}$, the Hamiltonian for $N$
electrons in $N$ wells is
\[
  \hat H(t)=\sum_{\sigma=1}^{N}\Bigl[-\tfrac12\nabla_\sigma^2
  +v^{\sigma}\bigl(\boldsymbol r_\sigma-\boldsymbol r_c^{\sigma};t\bigr)\Bigr]
  +\sum_{\sigma<\tau}\frac{\kappa}{|\boldsymbol r_\sigma-\boldsymbol r_\tau|}
  +\lambda(t)\,\hat O ,
\]
where each well is anharmonic along the register axis and harmonic in
the transverse direction,
\[
  v^{\sigma}(x,y)=\tfrac12(\omega^{\sigma})^2x^2
  +\tfrac13\beta^{\sigma}(\omega^{\sigma})^2x^4
  +\tfrac12(\omega_y^{\sigma})^2y^2 ,
\]
with frequencies of 5--15~GHz and well separations of a few
micrometers as in Refs.~[3,7].  This form applies when the barriers
suppress tunneling, so that the electrons are distinguishable; with a
smooth multi-well potential common to all electrons the wave function
is antisymmetric and the spin enters.  The last term is the external
probe.
We distinguish two cases.
\begin{itemize}
  \item \textbf{Time-independent probes.}  A static electric field, a
  static detuning of the well frequencies, or a perpendicular magnetic
  field.  They shift and mix the levels and define the operating
  points of the device: idle points, where the qubits are decoupled,
  and target points, where two wells exchange an excitation.  A
  time-independent probe that is switched on suddenly also defines
  the simplest dynamical problem, the evolution after a quench.
  \item \textbf{Time-dependent probes.}  Smooth ramps of the well
  parameters between operating points, which realize two-qubit gates
  such as $\sqrt{i\mathrm{SWAP}}$~[3,7], and resonant microwave fields
  acting on a single well, which realize single-qubit rotations and
  expose the leakage out of the qubit states.
\end{itemize}

In a PINN the wave function is a neural network that depends on the
coordinates $\boldsymbol R=(\boldsymbol r_1,\dots,\boldsymbol r_N)$ and
on time, with the initial state built into the ansatz, and the
network is trained by minimizing the residual of the Schr\"odinger
equation~[9,10],
\[
  \mathcal L(\theta)=\int_0^T\!dt\int\! d\boldsymbol R\;
  \bigl|\,i\,\partial_t\Psi_\theta(\boldsymbol R,t)
  -\hat H(t)\Psi_\theta(\boldsymbol R,t)\bigr|^2 ,
\]
sampled at collocation points in space and time.  For the
time-independent problem Ref.~[8] showed that a residual objective is
not variational: it gives no bound on the energy and has to be
followed by an energy minimization.  The time-dependent problem is
different.  For an ansatz with the exact initial state, the distance
to the exact solution is bounded by the accumulated residual,
\[
  \bigl\|\Psi_\theta(t)-\Psi(t)\bigr\|\le\int_0^t
  \bigl\|\bigl(i\,\partial_{t'}-\hat H(t')\bigr)\Psi_\theta(t')\bigr\|\,dt' ,
\]
so the training loss is at the same time an error
estimate~[11,12].  The same residual, minimized at each instant over
the time derivative of the parameters, is McLachlan's variational
principle~[13], which leads to time-dependent variational Monte
Carlo~[14--16]; its equations of motion contain the same matrix as the
stochastic-reconfiguration step of Ref.~[8].  A PINN that is trained
over a whole time interval at once and a neural-network state that is
stepped forward in time are thus two ways of minimizing one quantity,
with different strengths: the first has no accumulation of time-step
errors but must respect causality during training~[17] and resolve
rapidly oscillating phases, while the second is local in time but is
sensitive to noise and to the nodes of the wave function~[15].

The project addresses three questions:
\begin{enumerate}
  \item Can the PINN approach of Ref.~[8] be extended to multi-well
  potentials in two dimensions and to excited states, so that the
  qubit levels and their dependence on static probes are obtained
  with controlled accuracy?
  \item How accurately, and over how long times, can a PINN follow the
  dynamics of two interacting electrons under static and
  time-dependent probes, as measured against numerically exact grid
  calculations?
  \item How does the approach scale when more electrons and wells are
  added, and how does it compare with MCTDH(F) in accuracy, in cost
  and in what it can describe (two dimensions, bare Coulomb
  interaction, several electrons per well)?
\end{enumerate}

\section*{3. Objectives of the Thesis}

\subsection*{Objective 1: Theory and reproduction}
\begin{itemize}
  \item Study variational Monte Carlo and neural-network quantum
  states~[18--20], PINNs~[9,10,21], the time-dependent variational
  principle, and the MCTDH(F) method.
  \item Reproduce the two-electron results of Ref.~[8] for a single
  parabolic dot with the existing code.
  \item Set up grid-based reference solvers (discrete-variable
  representation, split-operator or Krylov propagation) for one
  electron in two dimensions and for two electrons in one and two
  dimensions.
\end{itemize}

\subsection*{Objective 2: Multi-well states and time-independent probes}
\begin{itemize}
  \item Adapt the ansatz to several wells: orbitals localized in the
  wells, the analytic electron--electron cusp, a neural correlator
  and backflow.  When tunneling between the wells is negligible the
  electrons can be treated as distinguishable, as in Ref.~[7]; the
  antisymmetric form is kept as the general case.
  \item Compute the ground state and the lowest excited states of two
  electrons in two wells, which define the two-qubit states, using
  residual pretraining followed by energy minimization and a method
  for excited states~[22].
  \item Study the response to static probes: level shifts in an
  electric field, detuning scans, avoided crossings, and the idle and
  target operating points.  Compare with exact diagonalization.
\end{itemize}

\subsection*{Objective 3: Time evolution of one and two electrons}
\begin{itemize}
  \item Develop a complex-valued, time-dependent ansatz
  $\Psi_\theta(\boldsymbol R,t)$ with time-dependent Jastrow and
  backflow terms~[14] and the initial state built in, and train it
  with the space--time residual, using causality-respecting weights
  and successive time windows.
  \item Test on one driven electron in an anharmonic well (Rabi
  oscillations and leakage).
  \item Two electrons, one per well: evolution after a quench under a
  time-independent Hamiltonian, exchange of an excitation between the
  wells at a target point, a ramped $\sqrt{i\mathrm{SWAP}}$ gate and
  a resonant drive.
  \item Monitor norm and energy conservation, the residual error
  bound, and the fidelity with the grid reference; determine the
  time window over which the PINN is reliable.
\end{itemize}

\subsection*{Objective 4: Global training against time stepping}
\begin{itemize}
  \item Implement time-dependent variational Monte Carlo for the same
  ansatz and compare with the space--time PINN in accuracy, stability
  and cost.
  \item Train networks that take the probe parameters (field
  strength, ramp and hold times) as additional inputs, so that one
  network describes a family of protocols.
\end{itemize}

\subsection*{Objective 5: More electrons and comparison with MCTDH(F)}
\begin{itemize}
  \item Extend the calculations to three and four wells with one
  electron per well in two dimensions, where exact grid calculations
  are no longer possible.
  \item Compare with the MCTDH(F) codes of O.~Leinonen, which are used
  as they are, at three levels: (i) the same one-dimensional
  Hamiltonian as in Ref.~[7], for a direct comparison of energies,
  populations and gate fidelities; (ii) the two-dimensional
  calculation with strong transverse confinement, which tests the
  effective one-dimensional interaction used in Ref.~[7]; (iii)
  genuinely two-dimensional situations (weak transverse confinement,
  wells that are not on a line), where the PINN results are
  predictions.
\end{itemize}

\subsection*{Objective 6: Possible extensions}
\begin{itemize}
  \item Several electrons per well (three spin-polarized electrons per
  well as in Ref.~[7]), where antisymmetry within each well is
  required.
  \item Lower barriers, where tunneling and the electron spin matter.
  \item A perpendicular magnetic field.
\end{itemize}

\section*{4. Methods and Tools}

\begin{itemize}
  \item \textbf{Neural networks:} the PyTorch code developed for
  Ref.~[8] (ansatz, residual training on stratified collocation
  points, stochastic reconfiguration) is the starting point; JAX is
  an alternative.  Production runs use GPUs.
  \item \textbf{Reference calculations:} grid solvers written in the
  project for small systems, and the MCTDH(F) codes of O.~Leinonen for
  larger ones.
  \item \textbf{Useful background:} quantum mechanics, computational
  physics with variational Monte Carlo (FYS4411), machine learning
  (FYS-STK4155) and many-body physics (FYS4480).
\end{itemize}
All codes and results will be documented in Jupyter notebooks so that
the thesis can be reproduced.

\section*{5. Work Plan and Milestones (One-Year Timeline)}

\subsection*{Semester 1 (Months 1--6)}

\textbf{Months 1--2: Theory and reproduction}
\begin{itemize}
    \item Literature study; reproduction of the two-electron results
    of Ref.~[8]; grid reference solvers.
\end{itemize}
\textbf{Deliverable:} Short report with reproduced results and
validated reference codes.

\bigskip

\textbf{Months 3--4: Multi-well states and static probes}
\begin{itemize}
    \item Multi-well ansatz; ground and excited states of two
    electrons in two wells; response to static probes; operating
    points.
\end{itemize}
\textbf{Deliverable:} Two-qubit level structure compared with exact
diagonalization.

\bigskip

\textbf{Months 5--6: Time-dependent PINN}
\begin{itemize}
    \item Time-dependent ansatz and space--time training; one driven
    electron; quench dynamics of two electrons.
\end{itemize}
\textbf{Deliverable:} Working time-dependent PINN validated against
grid calculations.

\subsection*{Semester 2 (Months 7--12)}

\textbf{Months 7--8: Time-dependent probes}
\begin{itemize}
    \item Ramped gates and resonant drives for two electrons;
    comparison of global training and time stepping; networks with
    the probe parameters as inputs.
\end{itemize}
\textbf{Deliverable:} Accuracy and cost study of the two training
strategies.

\bigskip

\textbf{Months 9--10: More electrons and MCTDH(F) comparison}
\begin{itemize}
    \item Three and four wells; comparison with MCTDH(F) in one
    dimension and in the quasi-one-dimensional limit; two-dimensional
    predictions.
\end{itemize}
\textbf{Deliverable:} Benchmark of PINN against MCTDH(F).

\bigskip

\textbf{Month 11: Extensions and consolidation}
\begin{itemize}
    \item One of the extensions of Objective~6, depending on progress;
    final production runs and figures.
\end{itemize}

\bigskip

\textbf{Month 12: Writing and final submission}
\begin{itemize}
    \item Completion of the thesis manuscript (methods and theory
    chapters are written continuously during the year) and
    preparation of the final presentation.
\end{itemize}
\textbf{Deliverable:} Completed thesis and defense presentation.

\section*{6. Expected Outcomes}

\begin{itemize}
    \item A documented PINN code for stationary states and time
    evolution of interacting electrons in two-dimensional multi-well
    potentials.
    \item Validated results for two electrons under static and
    time-dependent probes, with error estimates from the residual.
    \item A comparison of global space--time training and time
    stepping for neural-network wave functions.
    \item A benchmark against MCTDH(F) for several wells and a test of
    the one-dimensional effective interaction.
    \item A master thesis and, depending on the results, a
    contribution to a scientific article.
\end{itemize}

\section*{7. References}

\begin{enumerate}
    \item P.~M.~Platzman and M.~I.~Dykman, \textit{Quantum Computing with Electrons Floating on Liquid Helium}, Science \textbf{284}, 1967 (1999).
    \item N.~R.~Beysengulov \textit{et al.}, \textit{Coulomb Interaction-Driven Entanglement of Electrons on Helium}, PRX Quantum \textbf{5}, 030324 (2024).
    \item O.~Leinonen \textit{et al.}, \textit{Design and dynamics of two-qubit gates with motional states of electrons on helium}, Phys.\ Rev.\ A \textbf{113}, 032624 (2026).
    \item H.-D.~Meyer, U.~Manthe and L.~S.~Cederbaum, \textit{The multi-configurational time-dependent Hartree approach}, Chem.\ Phys.\ Lett.\ \textbf{165}, 73 (1990); M.~H.~Beck, A.~J\"ackle, G.~A.~Worth and H.-D.~Meyer, \textit{The multiconfiguration time-dependent Hartree (MCTDH) method: a highly efficient algorithm for propagating wavepackets}, Phys.\ Rep.\ \textbf{324}, 1 (2000).
    \item J.~Zanghellini, M.~Kitzler, T.~Brabec and A.~Scrinzi, \textit{Testing the multi-configuration time-dependent Hartree--Fock method}, J.\ Phys.\ B \textbf{37}, 763 (2004); J.~Caillat \textit{et al.}, \textit{Correlated multielectron systems in strong laser fields: A multiconfiguration time-dependent Hartree--Fock approach}, Phys.\ Rev.\ A \textbf{71}, 012712 (2005).
    \item A.~U.~J.~Lode, C.~L\'ev\^eque, L.~B.~Madsen, A.~I.~Streltsov and O.~E.~Alon, \textit{Colloquium: Multiconfigurational time-dependent Hartree approaches for indistinguishable particles}, Rev.\ Mod.\ Phys.\ \textbf{92}, 011001 (2020).
    \item O.~Leinonen, \textit{MCTDHF for simulations of multi-qubit systems}, in preparation; O.~Leinonen, PhD thesis, University of Oslo.
    \item A.~N.~Sekkelsten, B.~Fore, D.~H.~B.~Lima, M.~Hjorth-Jensen, J.~M.~Kim and A.~Lovato, \textit{Neural-network quantum states for two-dimensional quantum dots: benchmark energies and real-space structure from strong to weak confinement}, in preparation (2026); A.~N.~Sekkelsten, \textit{Physics-Informed Neural Quantum States for Quantum Dots}, MSc thesis, University of Oslo (2026).
    \item I.~E.~Lagaris, A.~Likas and D.~I.~Fotiadis, \textit{Artificial Neural Networks for Solving Ordinary and Partial Differential Equations}, IEEE Trans.\ Neural Networks \textbf{9}, 987 (1998).
    \item M.~Raissi, P.~Perdikaris and G.~E.~Karniadakis, \textit{Physics-informed neural networks: A deep learning framework for solving forward and inverse problems involving nonlinear partial differential equations}, J.\ Comput.\ Phys.\ \textbf{378}, 686 (2019).
    \item C.~Lubich, \textit{From Quantum to Classical Molecular Dynamics: Reduced Models and Numerical Analysis}, European Mathematical Society (2008).
    \item A.~Sinibaldi, D.~Hendry, F.~Vicentini and G.~Carleo, \textit{Time-dependent Neural Galerkin Method for Quantum Dynamics}, arXiv:2412.11778 (2024).
    \item A.~D.~McLachlan, \textit{A variational solution of the time-dependent Schr\"odinger equation}, Mol.\ Phys.\ \textbf{8}, 39 (1964).
    \item J.~Nys, G.~Pescia, A.~Sinibaldi and G.~Carleo, \textit{Ab-initio variational wave functions for the time-dependent many-electron Schr\"odinger equation}, Nat.\ Commun.\ \textbf{15}, 9404 (2024).
    \item A.~Sinibaldi, C.~Giuliani, G.~Carleo and F.~Vicentini, \textit{Unbiasing time-dependent Variational Monte Carlo by projected quantum evolution}, Quantum \textbf{7}, 1131 (2023).
    \item G.~Carleo, F.~Becca, M.~Schir\'o and M.~Fabrizio, \textit{Localization and Glassy Dynamics Of Many-Body Quantum Systems}, Sci.\ Rep.\ \textbf{2}, 243 (2012).
    \item S.~Wang, S.~Sankaran and P.~Perdikaris, \textit{Respecting causality for training physics-informed neural networks}, Comput.\ Methods Appl.\ Mech.\ Eng.\ \textbf{421}, 116813 (2024).
    \item G.~Carleo and M.~Troyer, \textit{Solving the quantum many-body problem with artificial neural networks}, Science \textbf{355}, 602 (2017).
    \item D.~Pfau, J.~S.~Spencer, A.~G.~D.~G.~Matthews and W.~M.~C.~Foulkes, \textit{Ab initio solution of the many-electron Schr\"odinger equation with deep neural networks}, Phys.\ Rev.\ Research \textbf{2}, 033429 (2020).
    \item D.~Luo and B.~K.~Clark, \textit{Backflow Transformations via Neural Networks for Quantum Many-Body Wave Functions}, Phys.\ Rev.\ Lett.\ \textbf{122}, 226401 (2019); E.~M.~Nordhagen, J.~M.~Kim, B.~Fore, A.~Lovato and M.~Hjorth-Jensen, \textit{Efficient solutions of fermionic systems using artificial neural networks}, Front.\ Phys.\ \textbf{11}, 1061580 (2023).
    \item T.~De~Ryck and S.~Mishra, \textit{Numerical analysis of physics-informed neural networks and related models in physics-informed machine learning}, Acta Numerica \textbf{33}, 633 (2024).
    \item D.~Pfau, S.~Axelrod, H.~Sutterud, I.~von~Glehn and J.~S.~Spencer, \textit{Accurate computation of quantum excited states with neural networks}, Science \textbf{385}, eadn0137 (2024).
\end{enumerate}

\end{document}
